% !TEX program = pdflatex
\documentclass[11pt]{article}

% ── Packages ──────────────────────────────────────────────────────────
\usepackage[margin=1in]{geometry}
\usepackage{amsmath,amssymb,amsthm}
\usepackage{booktabs}
\usepackage{graphicx}
\usepackage{hyperref}
\usepackage{cleveref}
\usepackage{xcolor}
\usepackage{caption}
\usepackage{subcaption}
\usepackage{multirow}
\usepackage{enumitem}
\usepackage[numbers]{natbib}
\bibliographystyle{plainnat}

\hypersetup{colorlinks=true,linkcolor=blue!60!black,citecolor=green!50!black,urlcolor=blue!70!black}

% ── Macros ────────────────────────────────────────────────────────────
\newcommand{\Einf}{E_\infty}
\newcommand{\Reval}{R_{\mathrm{eval}}}
\newcommand{\Rtrain}{R_{\mathrm{train}}}
\newcommand{\relLtwo}{\varepsilon_{\mathrm{rel}}}
\DeclareMathOperator*{\argmin}{arg\,min}

% ── Auto-populated result macros ──────────────────────────────────
\newcommand{\holdoutWinFraction}{15/33 (45\%)}
\newcommand{\holdoutMeanImprovement}{1.2\%}
\newcommand{\withinSliceMeanFrac}{0.46}
\newcommand{\withinSliceUnstableFrac}{22/27 (81\%)}
\newcommand{\totalRuns}{32{,}783}
\newcommand{\mainRuns}{31{,}370}
\newcommand{\ablationRuns}{1{,}413}

% ══════════════════════════════════════════════════════════════════════
\title{%
  Beyond Power Laws: Scaling-Law Diagnostics for Operator Learning%
}
\author{Lucas Perrier}
\date{\today \quad (Working Draft)}

\begin{document}
\maketitle

% ══════════════════════════════════════════════════════════════════════
\begin{abstract}
Empirical scaling laws have guided resource allocation in large-scale machine
learning, but their applicability to operator learning---where discretization
resolution introduces a third scaling axis alongside dataset size and model
capacity---remains poorly understood.
We conduct a comprehensive study across three 1D PDE benchmarks (Burgers,
Darcy, diffusion), three architectures (MLP, DeepONet, FNO), and
\totalRuns{} experimental configurations (training runs plus\ncross-resolution evaluations) spanning two orders of magnitude in each axis.
Rather than assuming power-law relationships \emph{a priori}, we fit six
candidate functional forms per scaling axis and use the corrected Akaike
Information Criterion (AICc) to let the data select the governing law.

We separate our findings by evidence strength.
The most robust result is that \emph{the direction of resolution scaling
is task-dependent}: higher resolution helps uniformly for
advection-dominated problems (Burgers: 35/35 slices improve),
is mixed for elliptic problems (Darcy), and actively
\emph{increases error} for smooth problems (Diffusion: 27/35
slices worsen for MLP and DeepONet).
Second, \emph{data scaling consistently shows bounded diminishing
returns}: across all 9 task--model combinations (8 on the densified
grid; Diffusion MLP on the sparser $n=5$ grid),
the AICc-selected winner belongs to a family with a finite error floor
$\Einf$, though the specific law identity (exponential vs.\
saturation vs.\ power) is statistically fragile and mostly
exploratory at per-slice sample sizes of 4--7.
Third, capacity scaling is largely negligible within our parameter range.

The framework's contribution is \emph{diagnostic}, not predictive:
in a held-out contest, multi-law fitting matches but does not
meaningfully outperform a power-law-only baseline (\holdoutWinFraction{}
wins, ${\sim}$\holdoutMeanImprovement{} mean improvement).
The value lies in identifying when resolution investment helps or
hurts, and in showing that operator-learning scaling should be treated
as a model-selection problem over law families rather than an
exponent-estimation exercise---though generalising beyond 1D PDEs
remains an open question.
\end{abstract}

% ══════════════════════════════════════════════════════════════════════
\section{Introduction}
\label{sec:intro}

Empirical scaling laws have become a central tool for understanding and
predicting the performance of large-scale machine learning
systems~\citep{kaplan2020scaling,hoffmann2022training}. In natural
language processing and computer vision, power-law relationships between
test loss and dataset size, parameter count, and compute budget guide
resource allocation decisions worth millions of dollars.

Scientific machine learning (SciML) faces an analogous resource
allocation problem with a twist: the practitioner must additionally
choose the \emph{discretization resolution} at which training data are
generated. Solving PDEs at higher resolution is expensive---often the
dominant cost---yet the relationship between resolution and downstream
operator-learning error is poorly characterized.

Neural operators~\citep{lu2021learning,li2021fourier} learn mappings
between infinite-dimensional function spaces and are often marketed as
``resolution-invariant.'' Two prominent architectures---DeepONet
\citep{lu2021learning} and the Fourier Neural Operator (FNO)
\citep{li2021fourier}---use fundamentally different strategies to
handle varying discretizations. Despite these architectural innovations,
it remains unclear \emph{when} resolution invariance translates into
practical benefits. Does doubling the training resolution halve the
error? Does it matter more than doubling the dataset? How does the
answer depend on the PDE?

A subtler question, largely absent from the literature, is whether
power laws are even the \emph{right functional form} for operator-learning
scaling. The power-law assumption $E \propto X^{-\alpha}$ originates
from statistical learning theory and information-theoretic arguments
in specific regimes. When the error--resource relationship follows
a different functional form---logarithmic, exponential,
saturating---forcing a power-law fit obscures the true scaling behavior
and leads to misleading exponent estimates.

\paragraph{Contributions.}
\begin{enumerate}
  \item We introduce \textbf{multi-law model selection for scaling
    analysis}: rather than assuming power laws, we fit six candidate
    functional forms (power, exponential, logarithmic, stretched
    exponential, saturation, linear) and use AICc to identify the
    governing law per axis, per task, per model. The \emph{identity}
    of the winning law becomes an interpretable diagnostic, subject to
    sample-size limitations we quantify.
  \item We report \textbf{\totalRuns{} training runs}
    (\mainRuns{} main grid plus \ablationRuns{} optimizer ablations)
    and cross-resolution evaluations across three PDE
    families, three architectures, and seven capacity levels
    (\Cref{tab:run-accounting}).
  \item We find that \textbf{data scaling on the densified grid
    consistently selects laws with a finite error floor}: the
    AICc-selected winner in all 9 task--model combinations (8 on the
    densified grid; Diffusion MLP on the sparser $n=5$ grid) belongs
    to a family with an $\Einf$ parameter. However, within-slice
    bootstrap analysis shows that the specific law identity is fragile
    at small $n$; the robust conclusion is that data scaling shows
    strong diminishing returns and is consistent with a bounded-decay
    regime.
  \item We show that \textbf{resolution scaling is non-universal}: it
    helps consistently only for Burgers, is mixed for Darcy, and
    \emph{actively hurts} for Diffusion (MLP, DeepONet). This is the
    most statistically robust finding in the study (4 of 5 robust
    law-selection cases occur on the resolution axis).
  \item We extract \textbf{estimated error floors}
    $\Einf$ from the best per-axis fits that suggest an
    architecture-dependent ranking (FNO $<$ MLP $<$ DeepONet),
    though these are extrapolated estimates conditional on the
    winning functional form, and cross-architecture comparisons
    are confounded by unequal effective parameter budgets at
    equivalent capacity labels.
\end{enumerate}


% ══════════════════════════════════════════════════════════════════════
\section{Related Work}
\label{sec:related}

\paragraph{Neural scaling laws.}
\citet{kaplan2020scaling} and \citet{hoffmann2022training} established
power-law scaling of cross-entropy loss with dataset tokens, parameter
count, and compute for autoregressive language models.
\citet{hestness2017deep} showed similar trends across vision and speech.
In scientific ML, \citet{perrier2025scaling} introduced scaling-law
diagnostics for physics-informed approaches (PINNs, HNNs) in
low-dimensional settings, finding that physics priors shift the error
floor rather than the scaling exponent. Our work extends this program to
infinite-dimensional operator learning, elevates resolution to a
scaling variable, and---critically---relaxes the power-law assumption
via multi-law model selection.

\paragraph{Methodological positioning.}
The dominant methodology in neural scaling laws assumes a power-law
functional form \emph{a priori} and focuses on estimating exponents
and proportionality constants~\citep{kaplan2020scaling,hoffmann2022training,hestness2017deep}.
Our work adopts a complementary perspective: rather than assuming
power laws, we treat the \emph{identity} of the best-fitting
functional form as itself informative. This family-selection-centric
view does not discard the traditional approach; instead, it includes
power laws as one of six candidates and explicitly tests when the
traditional baseline succeeds (Diffusion, where power laws dominate)
and when it fails (Burgers / Darcy data scaling, where
exponential and saturation families provide better fits).

\paragraph{Operator learning.}
DeepONet~\citep{lu2021learning} builds on the universal approximation
theorem for operators~\citep{chen1995universal}, factorizing the learned
operator into branch and trunk components. The Fourier Neural Operator
(FNO)~\citep{li2021fourier} parameterizes integral kernels via
truncated Fourier series, achieving state-of-the-art accuracy on
several PDE benchmarks. Recent extensions include U-shaped
architectures~\citep{wen2022ufno}, attention
mechanisms~\citep{hao2023gnot}, and physics-informed
variants~\citep{wang2021learning}.

\paragraph{Resolution and discretization in operator learning.}
A key selling point of neural operators is their ability to generalize
across resolutions. \citet{li2021fourier} demonstrated FNO's
zero-shot super-resolution capability. However, systematic studies of
how \emph{training} resolution affects accuracy---and whether more
resolution is always better---are lacking. Our work fills this gap, and
our finding that higher resolution can \emph{increase} error challenges
the prevailing wisdom.


% ══════════════════════════════════════════════════════════════════════
\section{Multi-Law Scaling Framework}
\label{sec:framework}

\subsection{Problem Setting}
We consider the supervised learning of solution operators
$\mathcal{G}^\dagger : \mathcal{U} \to \mathcal{V}$, where
$\mathcal{U}$ and $\mathcal{V}$ are Banach spaces of functions on
some domain $\Omega \subset \mathbb{R}^d$. Given $N$ input--output
pairs $\{(u_i, v_i)\}_{i=1}^N$ with
$v_i = \mathcal{G}^\dagger(u_i)$, a model $\mathcal{G}_\theta$ with
$D = |\theta|$ parameters is trained on data discretized at spatial
resolution $R$ (grid points per spatial dimension).

\subsection{Beyond Power Laws: Candidate Functional Forms}
\label{sec:candidates}

Prior work on neural scaling laws has assumed the power-law form
$E(X) = E_\infty + a\,X^{-\alpha}$, motivated by learning-theoretic
bounds. We relax this assumption and consider six candidate laws for
describing how error $E$ varies with a scaling variable
$X \in \{N, D, R\}$:

\begin{table}[h]
\centering
\caption{Candidate functional forms for scaling analysis.
$X$ denotes the scaling variable (dataset size $N$, parameter count $D$,
or resolution $R$). Each law is fit via nonlinear least squares with
physically motivated bounds.}
\label{tab:candidates}
\begin{tabular}{@{}llcl@{}}
\toprule
\textbf{Name} & \textbf{Form} & $k$ & \textbf{Regime} \\
\midrule
Power law     & $E_\infty + a\,X^{-\alpha}$ & 3
  & Algebraic convergence \\
Exponential   & $E_\infty + a\,e^{-\alpha X}$ & 3
  & Spectral convergence \\
Logarithmic   & $a - b\,\ln X$ & 2
  & Sub-algebraic (diminishing returns) \\
Stretched exp.\ & $E_\infty + a\,e^{-\alpha X^\beta}$ & 4
  & Interpolates power/exp \\
Saturation    & $E_\infty + a\,(1 + X/X_0)^{-\alpha}$ & 4
  & Sigmoidal with characteristic scale \\
Linear        & $a + b\,X$ & 2
  & Flat or monotone trend \\
\bottomrule
\end{tabular}
\end{table}

\noindent The number of free parameters $k$ varies from 2 to 4 across
candidates, making penalized model selection essential.

\paragraph{Physical motivation.}
Different functional forms arise from different approximation-theoretic
regimes:
\begin{itemize}
  \item \textbf{Power law} $E \propto N^{-\alpha}$: classical
    statistical convergence rates; $\alpha$ is governed by the effective
    dimensionality of the function space.
  \item \textbf{Logarithmic} $E \propto -\ln N$: arises when the
    approximation problem has high effective dimension, causing
    extremely slow convergence. Also characteristic of online learning
    in non-parametric settings.
  \item \textbf{Exponential} $E \propto e^{-\alpha N}$: spectral
    convergence rates for smooth functions; characteristic of methods
    that resolve the solution in an optimal basis.
  \item \textbf{Linear} with $b > 0$: indicates that increasing $X$
    \emph{hurts} performance---a pathological regime (e.g., overfitting
    to high-resolution noise).
\end{itemize}

\subsection{Model Selection via AICc}

For each slice (fixed values of the other two axes), we fit all six
candidates via nonlinear least squares
(\texttt{scipy.optimize.curve\_fit} with parameter bounds;
SciPy uses the trust-region reflective algorithm when bounds are
finite). We select
the winner using the corrected Akaike Information Criterion:
\begin{equation}
  \mathrm{AICc} = n \ln\!\bigl(\mathrm{RSS}/n\bigr) + 2k +
  \frac{2k(k+1)}{n - k - 1},
\end{equation}
where $n$ is the number of data points, $k$ the number of parameters,
and $\mathrm{RSS}$ the residual sum of squares. The AICc correction
is essential given our small per-slice sample sizes ($n = 4$--$7$).

\paragraph{Aggregation.}
For each task--model--axis triplet, we compute the Akaike weight
$w_i = \exp(-\Delta_i/2)\big/\sum_j\exp(-\Delta_j/2)$
for every slice, where $\Delta_i = \mathrm{AICc}_i - \mathrm{AICc}_{\min}$.
We report the mean Akaike weight distribution across slices (rather
than raw vote counts), the dominant law, its median~$R^2$, and the
median characteristic parameter ($\alpha$, $b$, etc.).
To assess selection stability, we run $B=500$ bootstrap replicates
(resampling slices with replacement) and report the probability with
which the dominant law retains its winning status.
Cases where the top-two mean Akaike weights differ by less than 15\%
are flagged as \emph{ambiguous}.
As a supplementary sensitivity check, we also compute the Bayesian
Information Criterion (BIC) for every slice and report AICc--BIC
concordance (fraction of slices where both criteria agree).
At our small sample sizes ($n = 4$--$7$), BIC's penalty is actually
\emph{weaker} than AICc's corrected penalty, so low concordance
should not be read as evidence against the AICc selection
(see \Cref{app:bic}).


% ══════════════════════════════════════════════════════════════════════
\section{Experimental Setup}
\label{sec:experiments}

\subsection{PDE Benchmarks}

\begin{table}[h]
\centering
\caption{PDE benchmark tasks. All tasks map an input function to an
output function on the same domain, discretized at resolution $R$.}
\label{tab:tasks}
\begin{tabular}{@{}llll@{}}
\toprule
\textbf{Task} & \textbf{PDE} & \textbf{Operator} & \textbf{Character} \\
\midrule
Burgers & $u_t + u u_x = \nu u_{xx}$ &
  $u_0(x) \mapsto u(x, T)$ & Advection-dominated \\
Darcy & $-\nabla \cdot (a \nabla u) = f$ &
  $a(x) \mapsto u(x)$ & Elliptic, multiscale \\
Diffusion & $u_t = \nu u_{xx}$ &
  $u_0(x) \mapsto u(x, T)$ & Diffusion-dominated \\
\bottomrule
\end{tabular}
\end{table}

\noindent
Burgers uses a spectral solver with integrating-factor RK4 and
$\frac{2}{3}$-rule dealiasing. Darcy uses a spectral Galerkin method.
Diffusion uses an exact spectral solution. All solvers are
resolution-parameterized, generating data at arbitrary $R$.

\subsection{Architectures}

\begin{table}[h]
\centering
\caption{Model architectures and their resolution handling.}
\label{tab:models}
\begin{tabular}{@{}llp{7cm}@{}}
\toprule
\textbf{Model} & \textbf{Type} & \textbf{Resolution handling} \\
\midrule
MLP & Discretization-tied &
  Flattened input of dimension $R$; parameter count $\propto R$ \\
DeepONet & Mesh-free &
  Branch encodes input at training $R$; trunk evaluates at arbitrary
  query points \\
FNO & Spectral &
  Global convolutions in Fourier space; truncation to $k_{\max}$ modes
  enables resolution transfer \\
\bottomrule
\end{tabular}
\end{table}

\subsection{Scaling Grid}
We sweep over five axes:
\begin{itemize}
  \item \textbf{Dataset size} $N \in \{50, 100, 500, 1{,}000, 5{,}000\}$
  \item \textbf{Resolution} $R \in \{32, 64, 128, 256\}$
  \item \textbf{Model capacity}: 7 levels (\textit{tiny} to
    \textit{xlarge}; $\sim$5K--500K parameters)
  \item \textbf{Data seeds}: $\{11, 22\}$
  \item \textbf{Training seeds}: $\{101, 202\}$ (Diffusion);
    $\{101, 202, 303\}$ (Burgers, Darcy)
\end{itemize}
The seed count varies by task: Diffusion uses $2 \times 2 = 4$
seed combinations per configuration, while Burgers and Darcy use
$2 \times 3 = 6$. The baseline grid (3~models $\times$ 7~capacities
$\times$ 5~data sizes $\times$ 4~resolutions) thus comprises
$1{,}680$ training runs for Diffusion and $2{,}520$ each for Burgers
and Darcy, totalling $6{,}720$ baseline training runs before
cross-resolution evaluation.
Each model is trained
for up to 5{,}000 epochs with early stopping (patience 200) using Adam
($\text{lr} = 10^{-3}$, weight decay $10^{-5}$). Results are averaged
over the seed combinations per configuration, yielding grouped
data points for analysis.

\paragraph{Extension grids.}
\Cref{tab:run-accounting} provides a complete accounting of all
experimental runs. Beyond the baseline, we ran
dense-$N$ extensions (6 intermediate $N$ values), wide-$R$
extensions (6 additional $R$ values), controlled-MLP experiments
(fixed parameter count), and optimizer ablations
(Adam with $2\times$ patience; SGD with cosine annealing).
The ``Runs'' column in \Cref{tab:run-accounting} counts each
(training configuration, evaluation resolution) pair as one entry,
including cross-resolution evaluations.
The grand total is \totalRuns{}.

\begin{table}[h]
\centering
\caption{Experimental run accounting. ``Baseline'' is the $5\times 4
\times 7$ data--resolution--capacity grid (4--6 seed combinations
per configuration depending on task), including cross-resolution
evaluations for DeepONet and FNO.
Extension grids densify specific axes or test
robustness.}
\label{tab:run-accounting}
\small
\begin{tabular}{@{}llrl@{}}
\toprule
\textbf{Component} & \textbf{Scope} & \textbf{Runs} &
\textbf{Purpose} \\
\midrule
Baseline grid & All tasks $\times$ 3 models & 15{,}120 &
  Core scaling analysis \\
Dense-$N$ extension & Burg/Darcy/Diff $\times$ DON/FNO & 11{,}578 &
  Data-axis law discrimination \\
Wide-$R$ extension & Burgers $\times$ 3 models & 2{,}352 &
  Resolution-axis law stability \\
MLP controlled & Burg+Diff $\times$ MLP-ctrl & 2{,}016 &
  Parameter-count confound \\
Dense-$N$ (MLP only) & Burg $\times$ MLP & 304 &
  MLP data-axis extension \\
\midrule
\emph{Main grid subtotal} & & \mainRuns{} & \\
\midrule
Adam $2\times$ patience & Burgers $\times$ 3 models & 756 &
  Optimizer sensitivity \\
SGD + cosine & Burgers $\times$ 3 models & 657 &
  Optimizer sensitivity \\
\midrule
\textbf{Grand total} & & \textbf{\totalRuns{}} & \\
\bottomrule
\end{tabular}
\end{table}

\subsection{Cross-Resolution Evaluation}
After training, each model is evaluated at all resolutions
$\Reval \in \{32, 64, 128, 256, 512\}$ except its training resolution,
yielding 13{,}440 additional evaluation metrics. For DeepONet, branch
inputs at non-training resolutions are linearly interpolated to
$\Rtrain$ before the forward pass. FNO natively handles arbitrary input
resolutions via Fourier truncation/zero-padding.

\subsection{Multi-Law Fitting Procedure}
For each of the three scaling axes, we extract per-slice data by
fixing the other two axes at each combination present in the grid.
This yields:
\begin{itemize}
  \item \textbf{Data axis ($E$ vs $N$)}: $7 \times 4 = 28$ slices per
    task--model pair (one per capacity$\times$resolution), each with
    $n = 5$ points.
  \item \textbf{Resolution axis ($E$ vs $R$)}: $7 \times 5 = 35$
    slices per pair, each with $n = 4$ points.
  \item \textbf{Capacity axis ($E$ vs $D$)}: $5 \times 4 = 20$
    slices per pair, each with $n = 7$ points.
\end{itemize}
All six candidate laws are fit to each slice and the AICc-best is
recorded. We identify the dominant law as the one with the highest
mean Akaike weight across slices (see \Cref{sec:candidates}), and
report slice-level vote counts as supporting detail alongside the
median $R^2$ and the median characteristic parameter.


% ══════════════════════════════════════════════════════════════════════
\section{Results}
\label{sec:results}

Of the \totalRuns{} runs (baseline grid plus dense-$N$,
wide-$R$, controlled-MLP, and ablation extensions),
582 diverged (all in Burgers, $<$3.3\% of that task); these are
excluded from fitting and grouped statistics.
We report relative $L^2$ error
$\relLtwo = \|v_{\mathrm{pred}} - v_{\mathrm{true}}\|_2 /
\|v_{\mathrm{true}}\|_2$ as the primary metric.

% ── 5.1: Data scaling ────────────────────────────────────────────────
\subsection{Data Scaling: Logarithmic or Bounded Decay}
\label{sec:data-scaling}

\Cref{tab:data-laws} summarizes the AICc-selected data-scaling law
for each task--model combination.

\begin{table}[h]
\centering
\caption{Data scaling: AICc-selected law, median~$R^2$,
mean Akaike weight of the winner ($\bar{w}_1$), and bootstrap
probability that the winner retains its status across $B\!=\!500$
replicates. ``Votes'' shows the top-2 law counts across all
slices per task--model pair. Burgers uses the widened grid
(63~slices per pair); Darcy and Diffusion use the base grid
(28~slices).}
\label{tab:data-laws}
\begin{tabular}{@{}llccccl@{}}
\toprule
\textbf{Task} & \textbf{Model} & \textbf{Winner} & $\tilde{R}^2$ &
$\bar{w}_1$ & \textbf{Boot.\%} & \textbf{Votes (top-2)} \\
\midrule
\multirow{3}{*}{Burgers}
  & MLP      & saturation  & 0.999 & 0.36 & 80\% & sat:25, pow:19 \\
  & DeepONet & \textbf{exponential} & 0.971 & 0.48 & 100\% & exp:36, pow:17 \\
  & FNO      & power       & 0.993 & 0.36 & 76\% & pow:25, log:20 \\
\midrule
\multirow{3}{*}{Darcy}
  & MLP      & \textbf{exponential} & 0.881 & 0.69 & 100\% & exp:28 \\
  & DeepONet & \textbf{exponential} & 0.886 & 0.36 & 94\% & exp:16, pow:9 \\
  & FNO      & stretched\_exp & 0.946 & 0.55 & 89\% & str:17, exp:10 \\
\midrule
\multirow{3}{*}{Diffusion}
  & MLP      & \textbf{power} & 0.998 & 0.68 & 100\% & pow:20, exp:5 \\
  & DeepONet & \textbf{power} & 0.909 & 0.47 & 100\% & pow:21, exp:6 \\
  & FNO      & \textbf{power} & 0.992 & 0.27 & 59\% & pow:11, str:10 \\
\bottomrule
\end{tabular}
\vspace{2pt}

\noindent{\footnotesize Burgers results include the widened resolution
grid ($R$ up to 512), yielding 63 slices per model. Darcy and
Diffusion use the base 4-resolution grid (28 slices per model).
Diffusion FNO remains the most ambiguous case: power wins with
only 0.27 mean weight and 59\% bootstrap probability. The within-slice
bootstrap (\Cref{tab:evidence-strength}) shows that most of these
law-identity selections are exploratory; the robust finding is
that all winners belong to bounded-decay families.}
\end{table}

\paragraph{Key finding: data scaling shows strong diminishing returns,
consistent with a bounded-decay regime.}
With the full analysis (including dense-$N$ and wide-$R$ extensions),
every task--model combination selects a law family that includes
a finite $\Einf$ term.
Logarithmic and linear forms---both lacking $\Einf$---are rejected
in every case.
On the sparse $n = 5$ grid used in earlier analysis, logarithmic
appeared dominant in 7 of 9 cases; with the denser grid
(\Cref{app:dense-n}), logarithmic is displaced entirely---
its apparent dominance was an
\textbf{artefact of sparse sampling} that could not resolve the
curvature near the error plateau.

\textbf{Important caveat:} while all winners share the $\Einf$
property, the within-slice bootstrap (\Cref{sec:within-slice-bootstrap})
shows that the specific law identity is \emph{fragile} at the
per-slice level. Only one data-axis case (Diffusion MLP: $\bar{w}_1
= 0.68$, within-slice dominant fraction $= 0.59$) meets both
robustness criteria. The reliable conclusion is not which specific
family wins (e.g., exponential vs.\ saturation vs.\ power---all
of these appear as winners in \Cref{tab:data-laws} depending on
task and model), but that (i)~data scaling follows a bounded-decay
pattern rather than unbounded logarithmic growth, and (ii)~error
decreases much more slowly than classical power-law rates suggest
for Burgers and Darcy.

\paragraph{Exception: Diffusion shows genuine power laws.}
All three Diffusion models exhibit power-law data scaling
(with floor): MLP ($\alpha \approx 1.1$, $R^2 = 0.998$),
FNO ($R^2 = 0.992$), and DeepONet ($R^2 = 0.909$). This is the only
task in our study where the classical power-law assumption holds
consistently. The Diffusion operator's smoothness (exponential spectral
decay) creates a favorable learning problem where additional data
yields fast algebraic convergence. DeepONet's lower $R^2$ and weaker
Akaike weight ($\bar{w}_1 = 0.47$) suggest it is less well-suited
to this favorable regime.

\begin{figure}[t]
\centering
\includegraphics[width=\textwidth]{../figures/multilaw_data_scaling.png}
\caption{Data scaling with AICc-selected fits overlaid on representative
slices. Burgers and Darcy include dense-$N$ extensions ($n \le 11$
data points); Diffusion uses the baseline $n=5$ grid.
Red: power law $\Einf + a N^{-\alpha}$. Purple: exponential.
Orange: saturation / stretched-exponential.
Titles show the law vote distribution across all slices.
Exponential and saturation dominate for Burgers and Darcy;
power laws win for Diffusion.}
\label{fig:data-scaling}
\end{figure}

\paragraph{Implications.}
Whether the data-axis form is logarithmic (as on the sparse grid) or
bounded-decay (as with denser sampling), the practical message is the
same: additional data yields diminishing returns far below the algebraic
rates assumed by classical theory. Bounded-decay forms further imply an
\emph{irreducible error floor} $\Einf > 0$ that no amount of training
data can breach without changing resolution, architecture, or
representation.


% ── 5.2: Resolution scaling ─────────────────────────────────────────
\subsection{Resolution Scaling: Task-Dependent, Non-Universal}
\label{sec:resolution-scaling}

\begin{table}[h]
\centering
\caption{Resolution scaling: AICc-selected law, the fraction of 35
slices where increasing $R$ reduces error (``R helps'') or increases
error (``R hurts''), and the bootstrap winner probability.
The specific law identity on this axis should be treated as
\emph{exploratory} rather than confirmatory, given the small
per-slice sample size ($n=4$).}
\label{tab:res-laws}
\begin{tabular}{@{}llcccc@{}}
\toprule
\textbf{Task} & \textbf{Model} & \textbf{Winner} & $\tilde{R}^2$ &
\textbf{R helps} & \textbf{R hurts} \\
\midrule
\multirow{3}{*}{Burgers}
  & MLP      & logarithmic & 0.887 & 35/35 & 0/35 \\
  & DeepONet & logarithmic & 0.862 & 35/35 & 0/35 \\
  & FNO      & linear      & 0.781 & 34/35 & 0/35 \\
\midrule
\multirow{3}{*}{Darcy}
  & MLP      & logarithmic & 0.511 & 25/35 & 9/35 \\
  & DeepONet & linear      & 0.518 & 20/35 & 14/35 \\
  & FNO      & logarithmic & 0.601 & 27/35 & 6/35 \\
\midrule
\multirow{3}{*}{Diffusion}
  & MLP      & logarithmic & 0.957 & 1/35 & \textbf{27/35} \\
  & DeepONet & logarithmic & 0.642 & 0/35 & \textbf{27/35} \\
  & FNO      & logarithmic & 0.317 & 0/35 & 2/35 \\
\bottomrule
\end{tabular}
\end{table}

\paragraph{Key finding: the \emph{direction} of resolution scaling
is task-dependent and the most statistically robust finding in
this study.}
With only $n=4$ resolution values per slice on the baseline grid,
AICc-based law selection has limited statistical power; the robust
finding is \emph{whether}
resolution helps, hurts, or is negligible---not which functional form
best describes the trend. The specific law labels (e.g.,
``logarithmic'' or ``linear'' in \Cref{tab:res-laws}) should be
treated as \emph{exploratory}: at $n=4$, even 2-parameter families
are only weakly distinguishable, and these labels should not be
interpreted as confident identifications of the true functional form.
A restricted-candidate-set ablation
(retaining only 2-parameter families: linear, logarithmic, power)
confirms that the directional conclusions survive (see
Supplementary~\Cref{tab:restricted-comparison}). For Burgers, where
the widened grid ($n = 10$) is available, the direction is confirmed
in 99.6\% of slices.

\textbf{Burgers (resolution helps uniformly).}
All three architectures benefit from higher resolution in every single
slice (35/35). The AICc-selected law is logarithmic
($E = a - b \ln R$ with $b > 0$, $b \approx 0.075$ for MLP and
DeepONet), though at $n=4$ this label is tentative.
The directional finding is what matters: the 100\% consistency
across all capacity levels and dataset sizes confirms that resolution
is genuinely beneficial for this advection-dominated PDE. This is
physically expected, as Burgers solutions develop steep gradients
and near-shock structures that require sufficient grid resolution
to represent.

\textbf{Darcy (resolution effect is mixed).}
Resolution helps in 20--27 of 35 slices depending on the model, but
actively hurts in 6--14 slices. The low $R^2$ values
(0.51--0.60) indicate that no simple functional form captures the
$E$-vs-$R$ relationship well. Darcy flow solutions are smooth enough
that $R = 32$ captures most of the spatial structure; additional
resolution provides marginal, noisy benefit that does not support a
reliable scaling law.

\textbf{Diffusion (resolution actively hurts).}
This is the most striking finding: for MLP and DeepONet, increasing
resolution \emph{increases} error in 27 of 35 slices. The
AICc-selected fit is logarithmic with $b < 0$ (negative coefficient),
consistent with error growing as $\ln R$---though again, the
specific functional form should be treated as tentative at $n=4$.
FNO is essentially flat ($R^2 = 0.32$, near the noise floor). Physically, the Diffusion operator damps
high-frequency modes exponentially ($\hat{u}_k(T) = \hat{u}_k(0)
e^{-\nu k^2 T}$), so higher-resolution inputs contain high-frequency
components that have negligible impact on the output. For
discretization-tied architectures (MLP) and branch-encoding
architectures (DeepONet), these extra input dimensions add noise
without providing useful signal, impairing generalization.

\begin{figure}[t]
\centering
\includegraphics[width=\textwidth]{../figures/multilaw_resolution_all_slices.png}
\caption{Resolution scaling: all slices overlaid with normalized error
$E(R)/E(R_{\min})$. Each line is one (capacity, $N$) slice. Dashed
gray line at 1.0 marks no change. \textbf{Burgers}: consistent
downward trend (resolution helps). \textbf{Darcy}: scattered around
1.0 (marginal effect). \textbf{Diffusion MLP/DeepONet}: upward trend
(resolution hurts). Annotations show the count of slices where
resolution helps.}
\label{fig:resolution-all-slices}
\end{figure}


% ── 5.3: Capacity scaling ───────────────────────────────────────────
\subsection{Capacity Scaling: Already in the Sufficient Regime}
\label{sec:capacity-scaling}

\begin{table}[h]
\centering
\caption{Capacity scaling: AICc-selected law. Most combinations select
linear (flat slope $b \approx 0$) or logarithmic with tiny coefficient,
indicating negligible capacity dependence in the range studied.}
\label{tab:cap-laws}
\begin{tabular}{@{}llcc@{}}
\toprule
\textbf{Task} & \textbf{Model} & \textbf{Winner} & $\tilde{R}^2$ \\
\midrule
\multirow{3}{*}{Burgers}
  & MLP      & logarithmic & 0.828 \\
  & DeepONet & linear      & 0.851 \\
  & FNO      & linear      & 0.812 \\
\midrule
\multirow{3}{*}{Darcy}
  & MLP      & linear      & 0.196 \\
  & DeepONet & linear      & 0.114 \\
  & FNO      & logarithmic & 0.437 \\
\midrule
\multirow{3}{*}{Diffusion}
  & MLP      & logarithmic & 0.147 \\
  & DeepONet & linear      & 0.869 \\
  & FNO      & \textbf{power} & 0.953 \\
\bottomrule
\end{tabular}
\end{table}

\paragraph{Key finding: capacity is not a bottleneck.}
In 8 of 9 task--model combinations, the winning capacity-scaling law is
either linear (flat, $b \approx 0$) or logarithmic with a tiny
coefficient ($b < 0.02$). This indicates that, within our parameter
range ($\sim$5K--500K), all architectures are already
capacity-sufficient for these 1D PDE tasks. Adding more parameters
does not significantly reduce error.

\paragraph{Exception: Diffusion FNO shows genuine capacity scaling.}
Diffusion FNO is the sole case where AICc selects a power law for
capacity scaling ($\alpha \approx 3.8$, $R^2 = 0.95$). This is
consistent with FNO's architecture: adding more Fourier modes and
larger hidden dimensions directly expands the spectral resolution of
the learned kernel. For a smooth operator like Diffusion, this
additional spectral capacity translates directly into lower error.

\begin{figure}[t]
\centering
\includegraphics[width=\textwidth]{../figures/multilaw_capacity_scaling.png}
\caption{Capacity scaling with AICc-selected fits. Most task--model
combinations show flat or weakly logarithmic trends. Diffusion FNO
(bottom center) is the exception, with genuine power-law capacity
scaling.}
\label{fig:capacity-scaling}
\end{figure}


% ── 5.4: Law selection heatmap ───────────────────────────────────────
\subsection{Law Selection as a Diagnostic}
\label{sec:law-heatmap}

\Cref{fig:heatmap} visualizes the fraction of slices selecting each
candidate law across all task--model--axis combinations. Several
patterns emerge:

\begin{enumerate}
  \item \textbf{Data axis}: exponential dominates for Darcy MLP
    (100\%) and Burgers DeepONet (57\%); power law dominates for
    Diffusion MLP (71\%) and DeepONet (75\%); saturation and
    stretched-exponential fill the remaining combinations.
    Logarithmic receives zero votes for models with the dense-$N$
    grid, and still captures 32\% for Burgers FNO.
    The split broadly correlates with PDE smoothness, but per-slice
    law identity remains mostly exploratory
    (\Cref{tab:evidence-strength}).
  \item \textbf{Resolution axis}: logarithmic and linear are the only
    serious contenders. Power law, exponential, and stretched exponential
    receive zero or negligible votes---resolution scaling is never
    well-described by classical convergence rates in our study.
  \item \textbf{Capacity axis}: the vote is fragmented (linear,
    logarithmic, power all appear), reflecting the fundamental
    weakness of the capacity signal in our parameter range.
\end{enumerate}

\begin{figure}[t]
\centering
\includegraphics[width=\textwidth]{../figures/multilaw_heatmap.png}
\caption{Law selection heatmap: fraction of slices (by AICc) selecting
each candidate law per task--model--axis combination. Darker cells
indicate stronger consensus. The data and resolution axes show clear
structure; the capacity axis is diffuse.}
\label{fig:heatmap}
\end{figure}


% ── 5.5: Irreducible error floors ───────────────────────────────────
\subsection{Irreducible Error Floors as Architecture-Dependent Estimates}
\label{sec:einf}

For laws that include an $\Einf$ parameter (power, exponential,
saturation, stretched exponential), the fitted floor provides an
\emph{estimate} of the error that would persist even with infinite
data---the architecture's approximation error for that task.
Important caveats apply: (i)~$\Einf$ depends on the assumed
functional form and is an extrapolation, not a direct measurement;
(ii)~laws without an $\Einf$ parameter (logarithmic, linear)
cannot contribute floor estimates, so $\Einf$ is available only
for the subset of slices where a floor-bearing law wins.
From the data-axis fits (where the signal is strongest), we extract:

\begin{table}[h]
\centering
\caption{Estimated irreducible error floors $\Einf$ from AICc-best
data-scaling fits (median across slices where a floor-bearing law
wins). These are extrapolated estimates, not direct measurements.
Lower is better.}
\label{tab:einf}
\begin{tabular}{@{}lccc@{}}
\toprule
& \textbf{MLP} & \textbf{DeepONet} & \textbf{FNO} \\
\midrule
Burgers   & 0.277 & 0.357 & \textbf{0.062} \\
Darcy     & 0.767 & 0.776 & \textbf{0.598} \\
Diffusion & 0.021 & 0.192 & \textbf{0.003} \\
\bottomrule
\end{tabular}
\end{table}

\paragraph{Key finding: $\Einf$ suggests an architecture-dependent
ranking, subject to both functional-form assumptions and
cross-architecture fairness caveats.}
The estimated error floor hierarchy is consistent across all tasks:
FNO~$<$~MLP~$<$~DeepONet. However, this ranking carries two
important limitations: (i)~$\Einf$ depends on the assumed functional
form and is an extrapolation, not a direct measurement; and
(ii)~``capacity level'' maps to different absolute parameter counts
across architectures (e.g., ``medium'' FNO has $\sim$287K parameters
vs.\ ``medium'' MLP at $\sim$34K), so part of FNO's advantage may
reflect its generally larger parameter budgets rather than purely
its inductive bias. The ranking should be interpreted as
\emph{suggestive within this experimental design}, not as a
controlled comparison at identical parameter counts.

The largest gap is on Diffusion, where FNO's $\Einf = 0.003$ is
63$\times$ lower than DeepONet's 0.192. This dramatic difference
arises because FNO's Fourier parameterization aligns with the
spectral structure of the Diffusion operator, while DeepONet's
pointwise trunk network cannot exploit this structure. However,
this comparison is not parameter-count controlled, and the floor
estimates would shift under a different winning law family; the
63$\times$ gap is conditional on the selected functional form.

On Darcy (the hardest task), all models have high floors ($> 0.6$),
indicating a fundamental approximation challenge that no architecture
in our study adequately addresses.

\begin{figure}[t]
\centering
\includegraphics[width=0.85\textwidth]{../figures/multilaw_e_inf.png}
\caption{Estimated irreducible error floors $\Einf$ from data-scaling
fits. FNO achieves the lowest estimated floor on every task. Darcy
has uniformly high floors, indicating an intrinsically hard
approximation problem. These are extrapolated estimates conditional
on the winning functional form.}
\label{fig:einf}
\end{figure}


% ── 5.6: Cross-resolution transfer ──────────────────────────────────
\subsection{Cross-Resolution Transfer}
\label{sec:cross-res}

We evaluate each trained model at $\Reval \in \{32, 64, 128, 256, 512\}$
to construct transfer matrices (\Cref{fig:cross-res}).

\paragraph{Burgers.}
FNO exhibits strong asymmetric transfer: training at $\Rtrain = 256$
yields significantly lower error at all evaluation resolutions compared
to training at $\Rtrain = 32$. This confirms that high-resolution
training data carries information about fine-scale features (steep
gradients) that the model leverages at any evaluation resolution.
DeepONet transfers more uniformly, consistent with its trunk-based
output representation.

\paragraph{Darcy and Diffusion.}
Both tasks show near-uniform transfer matrices, consistent with the
resolution-scaling finding: when the operator is smooth enough that
resolution barely matters, cross-resolution transfer is trivially
successful (but uninformative).

\begin{figure}[t]
\centering
\includegraphics[width=\textwidth]{../figures/fig_cross_res_per_task.png}
\caption{Cross-resolution transfer matrices. Each cell shows mean
$\relLtwo$ for training at $\Rtrain$ (row) and evaluating at
$\Reval$ (column). Burgers--FNO shows strong asymmetric transfer
(lower-right is best). Darcy and Diffusion are near-uniform.}
\label{fig:cross-res}
\end{figure}


% ── 5.7: Best absolute performance ──────────────────────────────────
\subsection{Best Absolute Performance}
\label{sec:best}

\begin{table}[h]
\centering
\caption{Best relative $L^2$ error per model per task. Note that the
best MLP configurations for Darcy and Diffusion prefer $R = 32$,
consistent with the finding that higher resolution is unhelpful or
harmful for these tasks.}
\label{tab:best}
\begin{tabular}{@{}llrrl@{}}
\toprule
\textbf{Task} & \textbf{Model} & \textbf{$\relLtwo$} & $D$ &
\textbf{Config} \\
\midrule
\multirow{3}{*}{Burgers}
  & MLP      & 0.199 & 263K & $N\!=\!5000, R\!=\!256$, xlarge \\
  & DeepONet & 0.267 & 461K & $N\!=\!5000, R\!=\!256$, xlarge \\
  & FNO      & \textbf{0.048} & 287K & $N\!=\!5000, R\!=\!256$, medium \\
\midrule
\multirow{3}{*}{Darcy}
  & MLP      & 0.668 & 82K & $N\!=\!5000, R\!=\!32$, large \\
  & DeepONet & 0.734 & 428K & $N\!=\!5000, R\!=\!128$, xlarge \\
  & FNO      & \textbf{0.525} & 287K & $N\!=\!1000, R\!=\!64$, medium \\
\midrule
\multirow{3}{*}{Diffusion}
  & MLP      & 0.018 & 82K & $N\!=\!5000, R\!=\!32$, large \\
  & DeepONet & 0.099 & 404K & $N\!=\!500, R\!=\!32$, xlarge \\
  & FNO      & \textbf{0.002} & 124K & $N\!=\!5000, R\!=\!256$, small-med \\
\bottomrule
\end{tabular}
\end{table}

FNO achieves the lowest error on all tasks within this experimental
design: 4.8\% on Burgers, 52.5\% on
Darcy, 0.2\% on Diffusion (noting that capacity labels map to
different absolute parameter counts across architectures). Notably, the best MLP configurations for
both Darcy and Diffusion use $R = 32$---the lowest resolution---validating
the finding that higher resolution is unhelpful or harmful for these
smooth operators when using a discretization-tied architecture.


% ── 5.8: Held-out prediction contest ───────────────────────────────
\subsection{Held-Out Prediction Contest}
\label{sec:holdout}

The preceding sections select scaling laws based on in-sample AICc
comparison. A natural objection is that multi-law selection
\emph{describes} the data better simply because it has more candidates
to choose from, without necessarily \emph{predicting} extrapolation
better. We convert this from a model-selection question into a
falsifiable prediction question.

\paragraph{Protocol.}
For each task--model--axis combination, we hold out ${\sim}20\%$ of
grid configurations (slices), fit (a)~a power-law-only model and
(b)~the AICc-best multi-law model on the remaining ${\sim}80\%$ of
slices (pooled), and evaluate out-of-sample RMSE on the held-out
slices. Results are averaged over all 33 task--model--axis
combinations (including MLP-controlled).

\paragraph{Results.}
\Cref{tab:holdout} summarizes the held-out prediction contest.
Multi-law selection achieves lower out-of-sample relative RMSE in
\holdoutWinFraction{} of combinations---essentially a coin flip---with
mean relative improvement of \holdoutMeanImprovement{} and
\emph{median improvement of zero}. In most cases the two approaches
perform comparably, with differences in the fourth decimal place.
The clearest multi-law advantage appears on the Darcy data axis,
where the fit for MLP improves from 0.113 to 0.101.

These results show that multi-law selection does not overfit to
in-sample points (its out-of-sample performance is comparable to
power-law-only), but it provides no meaningful \emph{predictive}
advantage either. \textbf{The framework's value lies entirely in its
diagnostic power}---the ability to identify which functional form
best describes each axis---not in superior extrapolation.
We caution against interpreting multi-law selection as an improved
prediction methodology; it is a descriptive and interpretive tool.

\begin{table}[h]
\centering
\caption{Held-out prediction contest: out-of-sample relative RMSE for
power-law-only vs.\ AICc-best multi-law. Bold indicates lower OOS
error.}
\label{tab:holdout}
\small
\begin{tabular}{@{}llccc@{}}
\toprule
\textbf{Task} & \textbf{Model} & \textbf{Axis} &
\textbf{PL OOS} & \textbf{ML OOS} \\
\midrule
Burgers & DeepONet   & data       & 0.192 & \textbf{0.192} \\
        &            & capacity   & \textbf{0.170} & 0.170 \\
        &            & resolution & 0.149 & \textbf{0.148} \\
        & FNO        & data       & 0.305 & \textbf{0.304} \\
        &            & capacity   & \textbf{0.338} & 0.338 \\
        &            & resolution & 0.267 & \textbf{0.262} \\
        & MLP        & data       & 0.147 & \textbf{0.147} \\
        &            & capacity   & 0.247 & \textbf{0.247} \\
        &            & resolution & 0.198 & \textbf{0.192} \\
        & MLP-ctrl   & data       & \textbf{0.382} & 0.384 \\
        &            & capacity   & \textbf{0.880} & 0.880 \\
        &            & resolution & 0.213 & \textbf{0.195} \\
\midrule
Darcy   & DeepONet   & data       & 0.137 & \textbf{0.134} \\
        &            & capacity   & \textbf{0.303} & 0.304 \\
        &            & resolution & \textbf{0.148} & 0.149 \\
        & FNO        & data       & 0.101 & \textbf{0.097} \\
        &            & capacity   & 0.264 & 0.264 \\
        &            & resolution & \textbf{0.327} & 0.327 \\
        & MLP        & data       & 0.113 & \textbf{0.101} \\
        &            & capacity   & \textbf{0.236} & 0.236 \\
        &            & resolution & 0.358 & \textbf{0.357} \\
\midrule
Diffusion & DeepONet & data       & 0.172 & 0.172 \\
          &          & capacity   & \textbf{0.211} & 0.216 \\
          &          & resolution & 0.322 & \textbf{0.322} \\
          & FNO      & data       & \textbf{0.118} & 0.118 \\
          &          & capacity   & 0.362 & 0.362 \\
          &          & resolution & \textbf{0.321} & 0.321 \\
          & MLP      & data       & \textbf{0.089} & 0.089 \\
          &          & capacity   & \textbf{0.428} & 0.439 \\
          &          & resolution & \textbf{0.344} & 0.344 \\
          & MLP-ctrl & data       & \textbf{0.989} & 0.990 \\
          &          & capacity   & \textbf{2.859} & 2.872 \\
          &          & resolution & 0.224 & \textbf{0.205} \\
\bottomrule
\end{tabular}
\end{table}


% ── 5.9: Within-slice bootstrap ─────────────────────────────────────
\subsection{Within-Slice Bootstrap Uncertainty}
\label{sec:within-slice-bootstrap}

The inter-slice bootstrap (\Cref{app:akaike}) measures consensus
\emph{across} slices that each have a fixed AICc winner. It does
not propagate fitting uncertainty from the $n$ data points within
each slice. Since $n = 4$--$7$ in our experiments, the per-slice
winner may be sensitive to individual data points.

\paragraph{Protocol.}
For each slice, we resample the $n$ data points with replacement
($B = 500$), refit all six candidate laws, and reselect the AICc
winner each time. This yields a per-slice law-selection distribution
that reflects fitting uncertainty directly.

\paragraph{Results.}
\Cref{tab:within-slice} reports aggregate within-slice winner
stability. The within-slice dominant fraction (mean fraction of
bootstrap replicates that select the same law winner) is
\withinSliceMeanFrac{} on average, lower than the inter-slice
bootstrap stability, confirming that individual slices carry
substantial law-selection uncertainty. In
\withinSliceUnstableFrac{} of axis-level combinations, the
within-slice dominant fraction falls below 50\%, indicating that
the per-slice data is insufficient to confidently distinguish
law families.

\begin{table}[h]
\centering
\caption{Within-slice bootstrap stability ($B=500$): mean fraction of
resamples selecting the dominant law. Values below 0.50 indicate the
per-slice data is insufficient to distinguish law families.}
\label{tab:within-slice}
\small
\begin{tabular}{@{}llccc@{}}
\toprule
\textbf{Task} & \textbf{Model} & \textbf{Data} &
\textbf{Capacity} & \textbf{Resolution} \\
\midrule
Burgers & DeepONet & 0.35 & 0.43 & \textbf{0.76} \\
        & FNO      & 0.47 & 0.29 & \textbf{0.52} \\
        & MLP      & 0.43 & 0.37 & \textbf{0.66} \\
\midrule
Darcy   & DeepONet & 0.33 & \textbf{0.54} & 0.41 \\
        & FNO      & 0.41 & \textbf{0.52} & 0.45 \\
        & MLP      & 0.35 & 0.45 & 0.44 \\
\midrule
Diffusion & DeepONet & 0.43 & 0.45 & \textbf{0.51} \\
          & FNO      & 0.45 & 0.42 & 0.45 \\
          & MLP      & \textbf{0.59} & 0.45 & \textbf{0.59} \\
\bottomrule
\end{tabular}
\end{table}

\paragraph{Implications.}
These results reinforce the paper's core message: at small $n$,
law-selection conclusions should be interpreted as indicative
rather than definitive. Of 27 task--model--axis combinations, only
4 satisfy both thresholds for robust selection (mean Akaike weight
$>0.5$ \emph{and} within-slice dominant fraction $>0.5$); 7 meet
one threshold (``suggestive''); and 16 meet neither (``exploratory'').
\Cref{tab:evidence-strength} classifies every finding accordingly.
The resolution axis yields the most robust selections
(3 of 4 robust cases), consistent with the directional resolution
finding being the study's strongest result.
Data-axis and capacity-axis law-identity claims should be treated
with appropriate caution; the reliable data-axis conclusion is
the \emph{direction} (bounded diminishing returns) rather than
the specific functional form.

\begin{table}[h]
\centering
\caption{Evidence-strength classification for every task--model--axis
finding. ``Robust'': mean Akaike weight $>0.5$ and within-slice
bootstrap fraction $>0.5$ (both thresholds met). ``Suggestive'':
one threshold met. ``Exploratory'': neither met. Law identity is
the AICc aggregate winner; parenthetical is the within-slice
bootstrap mode when it disagrees.}
\label{tab:evidence-strength}
\small
\begin{tabular}{@{}llllc@{}}
\toprule
\textbf{Task} & \textbf{Model} & \textbf{Axis} &
\textbf{AICc Winner} & \textbf{Strength} \\
\midrule
\multirow{9}{*}{Burgers}
  & DeepONet & data       & exponential & Exploratory \\
  & DeepONet & capacity   & logarithmic (bs: linear) & Exploratory \\
  & DeepONet & resolution & logarithmic & \textbf{Robust} \\
  & FNO      & data       & power       & Exploratory \\
  & FNO      & capacity   & logarithmic & Exploratory \\
  & FNO      & resolution & linear (bs: log) & Suggestive \\
  & MLP      & data       & saturation (bs: power) & Exploratory \\
  & MLP      & capacity   & logarithmic & Suggestive \\
  & MLP      & resolution & logarithmic & \textbf{Robust} \\
\midrule
\multirow{9}{*}{Darcy}
  & DeepONet & data       & exponential (bs: log) & Exploratory \\
  & DeepONet & capacity   & logarithmic (bs: linear) & Suggestive \\
  & DeepONet & resolution & linear (bs: log) & Exploratory \\
  & FNO      & data       & stretched exp (bs: log) & Suggestive \\
  & FNO      & capacity   & logarithmic & Suggestive \\
  & FNO      & resolution & logarithmic & Exploratory \\
  & MLP      & data       & exponential (bs: power) & Suggestive \\
  & MLP      & capacity   & logarithmic (bs: linear) & Exploratory \\
  & MLP      & resolution & logarithmic & Suggestive \\
\midrule
\multirow{9}{*}{Diffusion}
  & DeepONet & data       & power (bs: log) & Exploratory \\
  & DeepONet & capacity   & linear       & Exploratory \\
  & DeepONet & resolution & logarithmic & Suggestive \\
  & FNO      & data       & power       & Exploratory \\
  & FNO      & capacity   & power (bs: linear) & Exploratory \\
  & FNO      & resolution & logarithmic & Exploratory \\
  & MLP      & data       & power       & \textbf{Robust} \\
  & MLP      & capacity   & logarithmic (bs: linear) & Exploratory \\
  & MLP      & resolution & logarithmic & \textbf{Robust} \\
\bottomrule
\end{tabular}

\vspace{2pt}
\noindent{\footnotesize Robust: 4/27. Suggestive: 7/27. Exploratory:
16/27. The concentration of robust cases on the resolution axis
(3/4) reflects its directional stability; data-axis and capacity-axis
law identity is generally exploratory.}
\end{table}


% ══════════════════════════════════════════════════════════════════════
\section{Discussion}
\label{sec:discussion}

\subsection{The Scaling Law Type as a Diagnostic}

The central methodological contribution of this work is demonstrating
that the \emph{identity} of the best-fitting scaling law can provide
diagnostic information beyond the law's fitted parameters.
We separate our findings into three tiers of evidence strength:

\paragraph{Robust findings (supported by both Akaike weights and
within-slice bootstrap).}
\begin{itemize}
  \item \textbf{Resolution scaling direction}: resolution uniformly
    helps for Burgers (advection-dominated), is marginal for Darcy
    (elliptic), and actively hurts for Diffusion (smooth).
    The resolution axis yields 3 of 4 robust law-selection cases,
    with logarithmic scaling ($b > 0$ for Burgers, $b < 0$ for
    Diffusion) confirmed by both criteria.
  \item \textbf{Diffusion MLP data scaling} follows power-law decay
    with a floor ($\bar{w}_1 = 0.68$, bootstrap fraction $= 0.59$),
    the only data-axis case meeting both robustness criteria.
\end{itemize}

\paragraph{Suggestive findings (one robustness criterion met, or
consistent direction across multiple analyses).}
\begin{itemize}
  \item \textbf{Bounded-decay data scaling}: across all 9 task--model
    pairs (8 on the densified grid; Diffusion MLP on the sparser
    $n=5$ grid), the AICc winner belongs to a family
    with a finite $\Einf$ floor. While the specific law identity is
    statistically fragile (only 1/9 meets both robustness criteria),
    the \emph{direction}---bounded diminishing returns rather than
    unbounded logarithmic decay---is consistent across grid densities,
    tasks, and optimizer choices.
  \item \textbf{Flat capacity scaling}: in 8 of 9 cases, the
    capacity winner is linear or logarithmic with negligible
    coefficient, indicating that all architectures are
    capacity-sufficient in our parameter range.
\end{itemize}

\paragraph{Exploratory findings (neither robustness criterion met;
interpret with caution).}
\begin{itemize}
  \item The specific identity of the data-axis winner
    (e.g., exponential vs.\ saturation vs.\ power) for Burgers and
    Darcy. These distinctions require larger per-slice samples.
  \item Capacity-axis law identity for most combinations
    (15/27 exploratory overall).
\end{itemize}

\subsection{When Does Resolution Matter?}

Our results establish a clear trichotomy:
\begin{enumerate}
  \item \textbf{Resolution-critical} (Burgers): The PDE generates
    fine-scale features (steep gradients, near-shocks) that require
    resolution to represent. AICc selects logarithmic improvement
    ($b \approx 0.075$) and resolution helps in 100\% of slices.
    Practitioners should invest in both resolution and data.
  \item \textbf{Resolution-marginal} (Darcy): The PDE has moderate
    spatial complexity. Resolution helps in $\sim$70\% of slices but
    the effect is noisy ($R^2 \sim 0.5$). Pilot experiments at the
    target application should guide resolution allocation.
  \item \textbf{Resolution-harmful} (Diffusion): The PDE is
    spectrally smooth, making high-resolution inputs redundant or
    actively harmful. Practitioners should use the lowest resolution
    that resolves the relevant modes and invest in more data instead.
\end{enumerate}

\subsection{Architecture Recommendations}

\paragraph{Cross-architecture fairness caveat.}
Comparing architectures requires caution: ``capacity level'' maps to
different parameter counts across architectures (e.g., ``medium''
MLP at $R = 128$ has $\sim$34K parameters, while ``medium'' FNO has
$\sim$287K). The capacity grids were chosen to span a comparable
\emph{range} for each architecture, but absolute parameter counts
differ. FNO's consistent advantage may partly reflect its generally
larger parameter budgets at equivalent capacity labels. The
$\Einf$ ranking (FNO $<$ MLP $<$ DeepONet) is suggestive but should
not be interpreted as a controlled comparison at identical
parameter counts.

\paragraph{FNO.}
Within this experimental design, the Fourier Neural Operator achieves
the lowest absolute error and the lowest estimated $\Einf$ on every
task. Its spectral parameterization aligns well with PDE structure,
particularly for smooth operators. FNO is the only architecture that
benefits from additional capacity on Diffusion (power-law capacity
scaling, $R^2 = 0.95$) and shows the strongest cross-resolution
transfer on Burgers. However, FNO's generally larger parameter counts
at equivalent capacity labels mean that part of this advantage may
reflect budget rather than architecture. Its resolution scaling on
Burgers is weaker than MLP and DeepONet (linear rather than
logarithmic on the baseline grid), suggesting that FNO's spectral
smoothing may truncate some high-frequency information at higher
resolutions.

\paragraph{DeepONet.}
DeepONet shows the weakest performance in our study: highest estimated
$\Einf$ on all tasks, slow data scaling even on Diffusion (where
FNO/MLP achieve power laws), and no meaningful capacity scaling.
However, this comparison is not fully controlled: DeepONet's capacity
grid spans a different parameter range than FNO's, and the
branch--trunk architecture may require larger scales or different
hyperparameter tuning to realise its theoretical advantages.
This conclusion is limited to our experimental scope (1D PDEs,
$\leq$500K parameters); whether it persists for larger-scale
DeepONet variants, higher-dimensional problems, or problems with
more complex output queries remains open. The uniform
cross-resolution transfer is a structural advantage.

\paragraph{MLP.}
The simple MLP baseline is competitive with DeepONet on every task and
outperforms it on Burgers and Diffusion. Despite the limitation that
$D \propto R$, the MLP's direct parameterization avoids the bottlenecks
of the branch--trunk factorization. Its preference for $R = 32$ on
Diffusion is a natural consequence: lower resolution yields a smaller,
better-regularized network.

\subsection{Revisiting the Three-Axis Power Law}

For comparison with prior work, we also fit the additive power-law model
$E(N,D,R) = \Einf + aN^{-\alpha} + bD^{-\beta} + cR^{-\gamma}$
(\Cref{tab:3d-r2}). The nominally high $R^2$ values (0.79--0.95)
initially appear to validate this model. However, the multi-law
analysis reveals why this is misleading:

\begin{table}[h]
\centering
\caption{$R^2$ for the additive 3D power law. High values are an
artifact of the dominant data axis spanning most of the variance.}
\label{tab:3d-r2}
\begin{tabular}{@{}lccc@{}}
\toprule
& \textbf{MLP} & \textbf{DeepONet} & \textbf{FNO} \\
\midrule
Burgers   & 0.935 & 0.915 & 0.871 \\
Darcy     & 0.886 & 0.788 & 0.880 \\
Diffusion & 0.951 & 0.842 & 0.893 \\
\bottomrule
\end{tabular}
\end{table}

\begin{enumerate}
  \item The data axis dominates the total variance. Since data scaling
    \emph{is} genuine (whether logarithmic or algebraic), the 3D
    power law captures it as an $N^{-\alpha}$ term, inflating the
    overall $R^2$ even when the other axes contribute no signal.
  \item The $\beta$ estimates are unreliable: bootstrap 95\% CIs span
    orders of magnitude (e.g., FNO Burgers: $\beta \in [0.35, 3.91]$),
    confirming that capacity scaling is essentially undetermined.
  \item Values of $\gamma$ that hit the fitting bound ($\gamma = 5$)
    should be interpreted as ``no resolution effect,'' not as a
    meaningful exponent. On Diffusion, the additive power law masks the
    fact that resolution \emph{increases} error.
\end{enumerate}

\noindent The multi-law framework provides a more honest and
informative characterization by decomposing each axis independently
and letting the data choose the appropriate functional form.


\subsection{Relationship to Paper~1}

This work extends the scaling-law diagnostic framework of
\citet{perrier2025scaling} in three directions: (i)~from
finite-dimensional dynamics to infinite-dimensional operators,
(ii)~from two axes ($N$, $D$) to three axes ($N$, $D$, $R$), and
(iii)~from assumed power laws to multi-law model selection. The
central finding of Paper~1---that physics priors shift the error
floor---has an analog here: the architectural inductive bias of
FNO (spectral parameterization) shifts $\Einf$ by up to 63$\times$
compared to DeepONet, a far larger effect than any data or resolution
scaling.


% ══════════════════════════════════════════════════════════════════════
\section{Conclusion}
\label{sec:conclusion}

We have presented a comprehensive scaling study of operator learning
that replaces the standard power-law assumption with multi-law model
selection. Across \totalRuns{} training runs spanning baseline,
dense-$N$, wide-$R$, controlled-MLP, and optimizer ablation grids,
we separate our findings by evidence strength, from most to least
statistically reliable:

\paragraph{Robust conclusions (confirmed by both Akaike weights and
within-slice bootstrap): directional resolution scaling.}
\begin{enumerate}
  \item \textbf{Resolution is not universally beneficial.} Higher
    resolution helps for advection-dominated PDEs (Burgers: 35/35
    slices improve) but \emph{actively hurts} for smooth PDEs
    (Diffusion: 27/35 slices worsen for MLP and DeepONet).
    This is the study's most statistically robust finding
    (3 of 4 robust law-selection cases).
  \item \textbf{Diffusion MLP follows power-law data scaling with
    a floor} ($\alpha \approx 1.1$, $R^2 = 0.998$), the only
    data-axis case meeting both robustness criteria.
\end{enumerate}

\paragraph{Suggestive conclusions (one criterion met or consistent
direction across analyses): bounded data scaling and architecture
ranking.}
\begin{enumerate}[resume]
  \item \textbf{Data scaling shows bounded diminishing returns.}
    On the densified grid (8 of 9 task--model pairs; Diffusion MLP
    uses the sparser $n=5$ grid), every combination selects a
    law family with a finite $\Einf$ floor; however, the specific
    law identity is statistically fragile for 8 of 9 cases.
    The reliable inference is the \emph{direction} (bounded decay),
    not the specific family.
  \item \textbf{Capacity is not a bottleneck} for 1D PDE tasks at
    the scales we study (5K--500K parameters), except for FNO on
    Diffusion, which exhibits clear power-law capacity scaling.
  \item \textbf{Estimated error floors $\Einf$ rank architectures
    consistently}: FNO achieves the lowest estimated floor on every
    task, with the gap largest on Diffusion (63$\times$ better than
    DeepONet). These are extrapolated estimates conditional on
    both the winning functional form and the (uncontrolled)
    parameter budgets across architectures.
\end{enumerate}

\paragraph{Exploratory conclusions (neither criterion met; require
confirmation with larger samples): specific law-family identity.}
\begin{enumerate}[resume]
  \item \textbf{The specific law identity} (exponential vs.\
    saturation vs.\ stretched-exponential) on the data axis for
    Burgers and Darcy. The multi-law framework can distinguish
    bounded from unbounded decay, but finer distinctions within the
    bounded-decay family require more data points per slice.
\end{enumerate}

\paragraph{Limitations and future work.}
Our study is limited to 1D PDEs and a 1D Darcy formulation; the
``resolution can hurt'' finding should not be generalized to operator
learning broadly without at least one higher-dimensional benchmark.
In higher dimensions, resolution effects are amplified ($R^d$ cost
growth) and may shift the resolution-scaling balance. The baseline resolution grid
has only 4 training values ($R \in \{32, 64, 128, 256\}$); the
widened 10-value grid (\Cref{app:wide-res}) substantially improves
law-selection power but is currently available only for Burgers.
Future work should extend the wide-$R$ grid to all tasks.
Diffusion MLP baseline has only $n=5$ on the data axis (not yet
densified), so its data-scaling results rely on the sparser grid.

\paragraph{MLP resolution--capacity confound (a key limitation).}
For the MLP baseline, the parameter count scales linearly with input
resolution ($D \propto R$), creating a confound between the
resolution and capacity axes. The controlled-parameter experiment
(\Cref{app:mlp-controlled}) rules out this confound for the
``resolution hurts'' finding on Diffusion: degradation persists even
at fixed parameter count (\Cref{tab:mlp-controlled}). However,
the confound remains uncontrolled for Burgers and Darcy, where
resolution effects are directionally different.

The multi-law framework could be extended to fit the full 3D
relationship with axis-specific functional forms (e.g., logarithmic
in $N$, linear in $R$, flat in $D$) rather than forcing all axes to
share a single form. Finally, incorporating recent architectures
(Geo-FNO, GNOT, operator transformers) and 2D/3D PDE benchmarks
would test the generality of our conclusions.


% ══════════════════════════════════════════════════════════════════════
\section*{Acknowledgments}
Experiments were conducted on three RunPod GPU instances (NVIDIA A4000,
L4, and RTX 3090). Total experiment wall time was approximately 8 hours
across all pods.


% ══════════════════════════════════════════════════════════════════════
\begin{thebibliography}{20}

\bibitem[Chen and Chen(1995)]{chen1995universal}
T.~Chen and H.~Chen.
\newblock Universal approximation to nonlinear operators by neural networks
  with arbitrary activation functions and its application to dynamical systems.
\newblock \emph{IEEE Trans. on Neural Networks}, 6(4):911--917, 1995.

\bibitem[Hao et~al.(2023)]{hao2023gnot}
Z.~Hao, Z.~Wang, H.~Su, C.~Ying, Y.~Dong, S.~Liu, Z.~Cheng, J.~Song, and
  J.~Zhu.
\newblock {GNOT}: A general neural operator transformer for operator learning.
\newblock In \emph{ICML}, 2023.

\bibitem[Hestness et~al.(2017)]{hestness2017deep}
J.~Hestness, S.~Narang, N.~Ardalani, G.~Diamos, H.~Jun, H.~Kianinejad,
  M.~M.~A. Patwary, Y.~Yang, and Y.~Zhou.
\newblock Deep learning scaling is predictable, empirically.
\newblock \emph{arXiv:1712.00409}, 2017.

\bibitem[Hoffmann et~al.(2022)]{hoffmann2022training}
J.~Hoffmann, S.~Borgeaud, A.~Mensch, E.~Buchatskaya, T.~Cai, E.~Rutherford,
  D.~de~Las~Casas, L.~A. Hendricks, J.~Welbl, A.~Clark, T.~Hennigan,
  E.~Noland, K.~Millican, G.~van~den~Driessche, B.~Damoc, A.~Guy,
  S.~Osindero, K.~Simonyan, E.~Elsen, J.~W. Rae, O.~Vinyals, and
  L.~Sifre.
\newblock Training compute-optimal large language models.
\newblock In \emph{NeurIPS}, 2022.

\bibitem[Kaplan et~al.(2020)]{kaplan2020scaling}
J.~Kaplan, S.~McCandlish, T.~Henighan, T.~P. Brown, B.~Chess, R.~Child,
  S.~Gray, A.~Radford, J.~Wu, and D.~Amodei.
\newblock Scaling laws for neural language models.
\newblock \emph{arXiv:2001.08361}, 2020.

\bibitem[Li et~al.(2021)]{li2021fourier}
Z.~Li, N.~Kovachki, K.~Azizzadenesheli, B.~Liu, K.~Bhatt, A.~Stuart, and
  A.~Anandkumar.
\newblock Fourier neural operator for parametric partial differential equations.
\newblock In \emph{ICLR}, 2021.

\bibitem[Lu et~al.(2021)]{lu2021learning}
L.~Lu, P.~Jin, G.~Pang, Z.~Zhang, and G.~E. Karniadakis.
\newblock Learning nonlinear operators via {DeepONet} based on the universal
  approximation theorem of operators.
\newblock \emph{Nature Machine Intelligence}, 3(3):218--229, 2021.

\bibitem[Perrier(2025)]{perrier2025scaling}
L.~Perrier.
\newblock Scaling-law diagnostics for physics-informed machine learning.
\newblock Working paper, 2025.

\bibitem[Wang et~al.(2021)]{wang2021learning}
S.~Wang, H.~Wang, and P.~Perdikaris.
\newblock Learning the solution operator of parametric partial differential
  equations with physics-informed {DeepONets}.
\newblock \emph{Science Advances}, 7(40), 2021.

\bibitem[Wen et~al.(2022)]{wen2022ufno}
G.~Wen, Z.~Li, K.~Azizzadenesheli, A.~Anandkumar, and S.~M. Benson.
\newblock {U-FNO}---an enhanced {Fourier} neural operator-based deep-learning
  model for multiphase flow.
\newblock \emph{Advances in Water Resources}, 163:104180, 2022.

\end{thebibliography}


% ══════════════════════════════════════════════════════════════════════
\appendix
\section{Experimental Details}
\label{app:details}

\subsection{Capacity Grid}

\begin{table}[h]
\centering
\caption{MLP and DeepONet capacity grid (hidden layers $\times$ hidden
width). FNO uses a separate grid over modes, width, and number of
layers.}
\label{tab:capacity-grid}
\begin{tabular}{@{}lrrr@{}}
\toprule
\textbf{Level} & \textbf{MLP layers} & \textbf{Hidden width} &
\textbf{$\sim D$ (MLP, $R\!=\!128$)} \\
\midrule
tiny      & 2 & 32  & $\sim$5K \\
small     & 2 & 64  & $\sim$13K \\
small-med & 2 & 96  & $\sim$22K \\
medium    & 2 & 128 & $\sim$34K \\
med-large & 2 & 192 & $\sim$62K \\
large     & 2 & 256 & $\sim$82K \\
xlarge    & 3 & 256 & $\sim$148K \\
\bottomrule
\end{tabular}
\end{table}

\subsection{Training Configuration}
All models are trained with Adam (learning rate $10^{-3}$, weight decay
$10^{-5}$), batch size 64, and maximum 5{,}000 epochs with early
stopping (patience 200, monitoring validation relative $L^2$). Data is
split 80/10/10 into train/validation/test. Relative $L^2$ loss
$\|\hat{v} - v\|_2 / \|v\|_2$ is used for training.

\subsection{PDE Solver Details}

\paragraph{Burgers equation.}
$u_t + u u_x = \nu u_{xx}$ on $[0, 2\pi)$ with periodic boundary
conditions, $\nu = 0.01$, $T = 1.0$. Solved with integrating-factor
RK4 in Fourier space with $\frac{2}{3}$-rule dealiasing and
$\Delta t = 10^{-3}$. Random initial conditions are superpositions of
Fourier modes with random phases and amplitudes drawn from a $k^{-2}$
energy spectrum.

\paragraph{Darcy flow.}
$-\nabla \cdot (a(x) \nabla u(x)) = f(x)$ on $[0, 1]$ with Dirichlet
boundary conditions, $f(x) = 1$. The coefficient field $a(x)$ is drawn
from a log-normal random field. Solved via spectral Galerkin method.

\paragraph{Diffusion equation.}
$u_t = \nu u_{xx}$ on $[0, 2\pi)$ with periodic boundary conditions,
$\nu = 0.01$, $T = 1.0$. Exact spectral solution: each Fourier mode
$\hat{u}_k(t) = \hat{u}_k(0) e^{-\nu k^2 t}$. Same initial condition
distribution as Burgers.

\subsection{Multi-Law Fitting Details}

All six candidate laws are fit via \texttt{scipy.optimize.curve\_fit}
(Levenberg--Marquardt) with a maximum of 10{,}000 iterations. Bounds
are set to enforce physical plausibility: $\Einf \geq 0$, $a > 0$,
$0.001 \leq \alpha \leq 10$ for power laws. For the exponential law,
the scaling variable is normalized to $[0, 1]$ to avoid numerical
overflow; the fitted rate is then rescaled to original units.

AICc is computed as
$\mathrm{AICc} = n \ln(\mathrm{RSS}/n) + 2k + 2k(k+1)/(n-k-1)$,
where $n$ is the number of points per slice and $k$ is the number of
free parameters. The correction term is essential for our small sample
sizes ($n = 4$--$7$). When $n \leq k + 1$, AICc is set to $+\infty$,
effectively disqualifying overparameterized models.

BIC is computed as $\mathrm{BIC} = n \ln(\mathrm{RSS}/n) + k \ln n$.
For large $n$, BIC penalizes complexity more heavily than AICc;
however, at our small sample sizes ($n = 4$--$7$), an important
inversion occurs: BIC's penalty $k \ln n \approx 1.6k$ (at $n=5$)
is \emph{weaker} than AICc's corrected penalty
$2k + 2k(k+1)/(n-k-1)$, so BIC actually favors \emph{more} complex
models at these sample sizes---the opposite of its asymptotic
reputation. At $n = 5$, AICc is the more appropriate criterion
(it was designed specifically for small samples); BIC concordance
is reported for transparency, not as a robustness check in the
usual sense.


\subsection{Robustness: BIC Concordance}
\label{app:bic}

\Cref{tab:bic-concordance} reports the fraction of slices where AICc
and BIC select the same winning law. Because of the small-$n$ penalty
inversion described above, low concordance on axes with $n \leq 5$
should \emph{not} be interpreted as evidence against the AICc-selected
law; it reflects the well-understood divergence between the two
criteria at small sample sizes, where BIC's penalty is paradoxically
weaker.

\begin{table}[h]
\centering
\caption{AICc--BIC concordance rate: fraction of slices where both
criteria select the same winning law, per axis.}
\label{tab:bic-concordance}
\begin{tabular}{@{}llccc@{}}
\toprule
\textbf{Task} & \textbf{Model} & \textbf{Data} & \textbf{Resolution} &
\textbf{Capacity} \\
\midrule
Burgers   & MLP      & 4\%  & 3\%   & 90\% \\
Burgers   & DeepONet & 21\% & 3\%   & 85\% \\
Burgers   & FNO      & 11\% & 9\%   & 70\% \\
Darcy     & MLP      & 4\%  & 43\%  & 80\% \\
Darcy     & DeepONet & 11\% & 69\%  & 95\% \\
Darcy     & FNO      & 14\% & 34\%  & 65\% \\
Diffusion & MLP      & 4\%  & 97\%  & 45\% \\
Diffusion & DeepONet & 11\% & 100\% & 80\% \\
Diffusion & FNO      & 21\% & 74\%  & 90\% \\
\bottomrule
\end{tabular}

\vspace{4pt}
\noindent{\footnotesize The low concordance on the data axis
($n=5$) and resolution axis ($n=4$) reflects the small-$n$ penalty
inversion: at these sample sizes, BIC's complexity penalty is weaker
than AICc's corrected penalty, so BIC systematically favors more
complex models---the opposite of its large-$n$ behavior.
Capacity-axis concordance is generally high ($>$65\%), consistent
with the larger per-slice sample size ($n=7$) where the penalty
inversion is less severe. AICc remains the primary criterion
throughout; BIC concordance is reported for transparency.}
\end{table}


\subsection{Robustness: Restricted Candidate Set (Resolution Axis)}
\label{app:restricted}

Because the resolution axis has only $n=4$ data points per slice,
fitting 4-parameter models (stretched exponential, saturation) risks
overfitting. As a robustness check, we re-run law selection on the
resolution axis using only the three 2-parameter candidate families:
linear ($a + bX$, $k=2$), logarithmic ($a - b\ln X$, $k=2$), and
power ($\Einf + aX^{-\alpha}$, $k=3$, the simplest family with an
irreducible error floor).

\begin{table}[h]
\centering
\caption{Full (6-law) vs.\ restricted (3-law) winner on the resolution
axis. ``Direction'' is the qualitative finding (helps / hurts / flat).
Directional conclusions are robust to the candidate set.}
\label{tab:restricted-comparison}
\begin{tabular}{@{}llccc@{}}
\toprule
\textbf{Task} & \textbf{Model} & \textbf{Full winner} &
\textbf{Restricted winner} & \textbf{Direction robust?} \\
\midrule
Burgers   & MLP      & logarithmic & logarithmic & \checkmark \\
Burgers   & DeepONet & logarithmic & logarithmic & \checkmark \\
Burgers   & FNO      & linear      & linear      & \checkmark \\
Darcy     & MLP      & logarithmic & logarithmic & \checkmark \\
Darcy     & DeepONet & linear      & linear      & \checkmark \\
Darcy     & FNO      & logarithmic & logarithmic & \checkmark \\
Diffusion & MLP      & logarithmic & logarithmic & \checkmark \\
Diffusion & DeepONet & logarithmic & logarithmic & \checkmark \\
Diffusion & FNO      & logarithmic & power       & \checkmark \\
\bottomrule
\end{tabular}
\end{table}

\noindent All directional conclusions (Burgers: helps; Darcy: mixed;
Diffusion: hurts for MLP/DeepONet, flat for FNO) are preserved when
the candidate set is restricted, supporting the claim that the
resolution findings are robust even though the specific law identity
is exploratory.


\subsection{Akaike Weights and Bootstrap Stability}
\label{app:akaike}

\Cref{fig:akaike-weights} shows the mean Akaike weight distribution
per task--model--axis combination. Hatched bars indicate cases where
the top-two mean Akaike weights differ by less than 15\%
(``ambiguous''). \Cref{fig:bootstrap} shows the bootstrap winner
probability for the dominant law.

\begin{figure}[h]
\centering
\includegraphics[width=\textwidth]{../figures/multilaw_akaike_weights.png}
\caption{Mean Akaike weights across slices. Hatched bars indicate
ambiguous cases where the top-two laws are within 15\% of each other.}
\label{fig:akaike-weights}
\end{figure}

\begin{figure}[h]
\centering
\includegraphics[width=\textwidth]{../figures/multilaw_bootstrap_stability.png}
\caption{Bootstrap law-selection stability ($B=500$). Green bars
($>$80\%) indicate confident selection; orange (60--80\%) moderate;
red ($<$60\%) ambiguous.}
\label{fig:bootstrap}
\end{figure}


% ══════════════════════════════════════════════════════════════════════
% TIER 3: Extended Experiments (Appendix)
% ══════════════════════════════════════════════════════════════════════

\section{Extended Experiments}
\label{app:tier3}

The following experiments extend the main study with new training runs
designed to address specific limitations identified in the core analysis.

\subsection{Densified Data Axis}
\label{app:dense-n}

The data axis in the main experiments uses $n = 5$ values of $N$
($\{50, 100, 500, 1000, 5000\}$). At this sample size,
AICc-based law selection can be brittle (\Cref{sec:within-slice-bootstrap}).
We add six intermediate values ($N \in \{150, 200, 300, 750, 2000, 3000\}$),
yielding $n = 11$ data points per slice---enough for more reliable
discrimination between logarithmic and power-law decay.

\paragraph{Key questions.}
\begin{enumerate}
  \item Does the ``logarithmic wins'' finding on the data axis
    stabilize or shift towards power-law with denser sampling?
  \item Does BIC concordance improve at $n = 11$ (where the small-$n$
    penalty inversion from \Cref{app:bic} becomes less severe)?
  \item Do within-slice bootstrap dominant fractions increase
    substantially with 11 data points?
\end{enumerate}

\paragraph{Results.}
Extending from $n = 5$ to $n = 11$ data-axis sample sizes
dramatically shifts law selection on the data axis.
With the original sparse grid ($N \in \{50, 100, 500, 1000, 5000\}$),
logarithmic was the dominant winner across Burgers
(DeepONet: 79\%, FNO: 52\%, MLP: 75\%) and Darcy (MLP: 89\%).
With the dense grid ($N \in \{50, 100, 150, 200, 300, 500, 750,
1000, 2000, 3000, 5000\}$), logarithmic is displaced: exponential
now dominates Burgers DeepONet (57\%) and both Darcy models
(DeepONet: 54\%, MLP: 96\%); power or saturation emerge for
Burgers FNO (power, 40\%) and Burgers MLP (saturation, 40\%).

All the new winners---exponential, saturation, stretched
exponential---encode a finite irreducible-error floor that the
logarithmic form lacks. With only five sample sizes, the
within-regime curvature is too coarsely sampled for AICc to separate
a log curve from a bounded-decay curve, biasing selection toward the
simpler logarithmic form. Intermediate points (especially $N = 150$,
200, 300, 750) resolve the curvature near the floor and tip
selection toward bounded families.

\emph{Takeaway:} With only five data-axis sample sizes, AICc
favors logarithmic; with the denser 11-point grid,
data-axis scaling is better described by families
with an irreducible-error floor, consistent with finite-resolution
approximation limits. The main-text \Cref{tab:data-laws} reports
results on the combined grid.
Dense-$N$ extensions are complete for all task--model pairs except
Diffusion MLP baseline (which retains only $n = 5$ data-axis
sample sizes). Darcy DeepONet and FNO now use the full $n = 31$
dense grid; Diffusion DeepONet and FNO use an $n = 25$ grid.

\subsection{MLP Fixed-Parameter-Count Experiment}
\label{app:mlp-controlled}

For the standard MLP, the parameter count $D$ scales linearly with
resolution $R$ (since $D = 2Rh + h^2 + 2h + R$ for hidden width $h$).
This creates a confound: increasing $R$ simultaneously increases input
dimensionality and model capacity. The \texttt{mlp\_controlled} model
variant adjusts the hidden width at each resolution so that $D$ remains
approximately constant (matched to the standard MLP at $R = 64$).

\paragraph{Protocol.}
Train \texttt{mlp\_controlled} on the Diffusion task (where ``resolution
hurts'' is strongest) across all capacity levels and
$R \in \{32, 64, 128, 256\}$. Compare against the standard
\texttt{mlp\_baseline} at the same configurations.

\paragraph{Interpretation.}
\begin{itemize}
  \item \textbf{If ``resolution hurts'' persists} with
    \texttt{mlp\_controlled} $\Rightarrow$ genuine resolution effect
    (extraneous high-frequency content harms learning).
  \item \textbf{If ``resolution hurts'' disappears} $\Rightarrow$ the
    confound was real (increasing $D$ with $R$ was the bottleneck).
\end{itemize}

\paragraph{Results.}
\Cref{tab:mlp-controlled} confirms that ``resolution hurts'' persists
even when parameter count is held approximately constant.
Across all seven capacity levels on the Diffusion task, the controlled
MLP achieves its lowest test error at $R = 32$ and progressively
degrades as $R$ increases to 256 --- the same qualitative pattern as
the standard MLP baseline. For example, at \texttt{medium} capacity the
controlled model holds $D \approx 33{,}100$ parameters regardless of
$R$, yet error rises from 0.021 ($R = 32$) to 0.026 ($R = 256$). The
baseline MLP shows the same direction (0.082 $\to$ 0.098)
despite its parameter count growing from 24{,}864 to 82{,}432.

The effect is most dramatic at small model sizes (\texttt{tiny}:
0.023 $\to$ 0.295), where the controlled model must compress all
information into very few parameters.
These results rule out the parameter-count confound:
\textbf{``resolution hurts'' is a genuine resolution effect on Diffusion},
consistent with the interpretation that the model wastes capacity
fitting output-irrelevant high-frequency content present at higher $R$.

\begin{table}[h]
\centering
\caption{Test relative $L^2$ error on Diffusion for the parameter-controlled MLP
(\texttt{mlp\_controlled}, top) vs.\ standard MLP baseline (bottom). Both show
monotonically increasing error with $R$, ruling out the parameter-count confound.}
\label{tab:mlp-controlled}
\small
\begin{tabular}{lcccc}
\toprule
Capacity & $R{=}32$ & $R{=}64$ & $R{=}128$ & $R{=}256$ \\
\midrule
\multicolumn{5}{l}{\emph{MLP Controlled (fixed parameter count)}} \\
tiny      & 0.023 & 0.038 & 0.140 & 0.295 \\
small     & 0.024 & 0.023 & 0.025 & 0.111 \\
small-med & 0.023 & 0.026 & 0.027 & 0.033 \\
medium    & 0.021 & 0.027 & 0.027 & 0.026 \\
med-large & 0.020 & 0.025 & 0.033 & 0.031 \\
large     & 0.018 & 0.024 & 0.031 & 0.034 \\
xlarge    & 0.027 & 0.032 & 0.039 & 0.047 \\
\midrule
\multicolumn{5}{l}{\emph{MLP Baseline (standard)}} \\
tiny      & 0.118 & 0.118 & 0.125 & 0.132 \\
small     & 0.096 & 0.101 & 0.105 & 0.103 \\
small-med & 0.088 & 0.094 & 0.098 & 0.100 \\
medium    & 0.082 & 0.091 & 0.099 & 0.098 \\
med-large & 0.074 & 0.086 & 0.095 & 0.101 \\
large     & 0.070 & 0.081 & 0.093 & 0.100 \\
xlarge    & 0.102 & 0.117 & 0.127 & 0.133 \\
\bottomrule
\end{tabular}
\end{table}

\subsection{Widened Resolution Grid}
\label{app:wide-res}

The main study uses $n = 4$ resolution values
($R \in \{32, 64, 128, 256\}$), making the law identity on the
resolution axis exploratory. We extend to 10 values:
$R \in \{16, 32, 48, 64, 96, 128, 192, 256, 384, 512\}$.

\paragraph{Key questions.}
\begin{enumerate}
  \item Does the Burgers / Darcy / Diffusion resolution trichotomy
    survive with a wider range?
  \item Are there regime changes (e.g., resolution helps at first then
    saturates)?
  \item Does the specific law identity (e.g., logarithmic vs.\ linear)
    become more robust at $n = 10$?
\end{enumerate}

\paragraph{Results.}
For Burgers (the only task with the full 10-R grid so far),
the widened grid reveals a pronounced shift in law identity.
Compared with the original 4-R grid:
\begin{itemize}
  \item \textbf{DeepONet:} logarithmic drops from 87\% $\to$ 51\%;
    linear rises from 12\% $\to$ 30\%.
  \item \textbf{FNO:} logarithmic drops from 58\% $\to$ 19\%;
    \textbf{linear emerges as the dominant winner} at 81\%.
  \item \textbf{MLP:} logarithmic drops from 92\% $\to$ 53\%;
    linear rises from 8\% $\to$ 39\%.
\end{itemize}
The directional finding is stable: resolution helps in 99.6\% of
Burgers slices (230/231) regardless of model or capacity.
The shift toward linear scaling for FNO is physically interpretable:
FNO operates in spectral space, so adding resolution linearly
increases the number of retained Fourier modes.

\emph{Takeaway:} The specific law family on the resolution axis
is sensitive to grid density. With 10 resolution values,
FNO resolution scaling is clearly linear, while MLP and DeepONet
retain logarithmic as a slight plurality winner but with significant
linear competition. Wide-$R$ data currently available for Burgers
only; Darcy and Diffusion use the extended grids from the main
baseline (DeepONet/FNO: $n = 20$; MLP: $n = 4$).

\subsection{Frequency-Filtered Diffusion}
\label{app:filtered}

To test the causal mechanism behind ``resolution hurts'' on Diffusion,
we apply a low-pass filter to the input initial conditions before
training. The filter retains only the first $k_{\max}$ Fourier modes
(default: $k_{\max} = 10$, matching the number of IC generation modes).
At high $R$, the standard task includes resolved but output-irrelevant
high-frequency content; the filtered variant removes it.

\paragraph{Protocol.}
Train all three architectures on \texttt{diffusion\_filtered} at
$R \in \{32, 64, 128, 256, 512\}$ with
$N \in \{500, 1000, 2000, 5000\}$.

\paragraph{Interpretation.}
\begin{itemize}
  \item \textbf{If filtering fixes the degradation at high $R$}
    $\Rightarrow$ strong support for ``extraneous frequencies hurt.''
  \item \textbf{If degradation persists} $\Rightarrow$ the mechanism
    is not purely spectral (e.g., optimisation difficulty scales with
    $R$ regardless of content).
\end{itemize}

\paragraph{Results.}
\texttt{[Placeholder: populate after running Tier 3.4 experiments
with configs/diffusion\_filtered.yaml]}

\subsection{DeepONet Interpolation Ablation}
\label{app:deeponet-ablation}

In the main experiments, cross-resolution evaluation of DeepONet
uses linear interpolation to map $R_{\mathrm{eval}}$-resolution
inputs to the $R_{\mathrm{train}}$-resolution branch input. This
could introduce interpolation artifacts, especially when
$R_{\mathrm{eval}} \gg R_{\mathrm{train}}$.

We compare:
\begin{enumerate}
  \item[(a)] \textbf{Interpolation} (standard): input interpolated to
    $R_{\mathrm{train}}$ sensor grid, trunk evaluates on
    $R_{\mathrm{eval}}$ query grid.
  \item[(b)] \textbf{Subsampling}: instead of interpolation, take
    evenly spaced $R_{\mathrm{train}}$ points from the
    $R_{\mathrm{eval}}$-resolution input. No smoothing.
\end{enumerate}

\paragraph{Results.}
We evaluated 288 cross-resolution comparisons across Darcy and
Diffusion tasks (3 capacities $\times$ 4 training resolutions $\times$
4--5 evaluation resolutions $\times$ 4 dataset sizes).

\textbf{When $R_{\mathrm{eval}} > R_{\mathrm{train}}$ (upsampling):}
interpolation and subsampling perform comparably (mean relative
$L^2$: 0.523 vs.\ 0.520), with subsampling winning 68\% of the 180
comparisons (by a slim margin).
\textbf{When $R_{\mathrm{eval}} < R_{\mathrm{train}}$ (downsampling):}
interpolation is overwhelmingly better (98\% of 108 cases), because
subsampling applied to coarser inputs introduces severe artifacts
(mean relative $L^2$: 0.531 vs.\ 1.068).

\paragraph{Conclusion.}
The standard interpolation strategy used in the main experiments is
robust: it never catastrophically degrades and is clearly preferable
for downsampling. For upsampling, the two methods are comparable,
so cross-resolution transfer conclusions in the main paper are
\textbf{not driven by interpolation artifacts}.

\subsection{Early Stopping and Optimizer Sensitivity}
\label{app:optimizer}

All main experiments use Adam with learning rate $10^{-3}$ and
early-stopping patience of 200 epochs. We test sensitivity by
rerunning a subset (Burgers, all 3 models, all capacities) with:
\begin{itemize}
  \item \textbf{Variant A}: Adam with 2$\times$ patience (400 epochs)
  \item \textbf{Variant B}: SGD with momentum 0.9 + cosine annealing schedule
\end{itemize}

\paragraph{Key question.}
Do law-selection winners change when the optimizer or convergence
criterion shifts? If they do, the law selection is training-sensitive
rather than reflecting a property of the task--architecture pair.

\paragraph{Results.}
We analyse Burgers at medium capacity (756 runs for
Adam~$2\times$; 657/756 complete for SGD~cosine).

\textbf{Resolution axis:} Law selection is fully stable across
optimizers. Logarithmic wins in 88\% of baseline slices,
90\% for Adam~$2\times$, and 100\% for SGD~cosine. This is
the strongest invariance result in the study.

\textbf{Data axis:} Law identity is more sensitive.
DeepONet shifts from power (baseline) to exponential (both variants).
FNO stays power under Adam~$2\times$ but shifts to logarithmic
under SGD. MLP shifts from saturation (baseline) to power (both
variants). However, all winners belong to bounded-decay families
(exponential, power, saturation), so the qualitative conclusion ---
data scaling is bounded, not unbounded-logarithmic --- is
optimizer-invariant.

\emph{Takeaway:} The resolution-axis law identity (logarithmic)
is robust to optimizer and convergence criterion. The data-axis
law \emph{family} (bounded decay vs.\ unbounded) is also stable,
though the specific law within that family is optimizer-sensitive.
Pending: diffusion optimizer ablation results.


\end{document}
